\documentclass[10pt, aspectratio=169]{beamer}
\usepackage{../beamer_theme_408}

\title{408考研零基础 --- 计算机组成原理}
\subtitle{2-8 Cache高速缓存}
\author{}
\date{}

\begin{document}


\begin{frame}{本节目标}
    \begin{enumerate}
        \item 理解Cache的工作原理——解决CPU-主存速度不匹配
        \item 掌握时间局部性和空间局部性
        \item 了解Cache的基本结构（标记、数据、有效位）
        \item 掌握Cache的读写策略（写直达、写回）
    \end{enumerate}
\end{frame}

\begin{frame}[t]{为什么需要Cache？}
    \begin{block}{CPU-主存速度不匹配}
        \begin{itemize}
            \item CPU速度不断提高（主频GHz级）
            \item 主存DRAM速度提升缓慢（几十ns）
            \item 差距可达\textbf{几个数量级}
        \end{itemize}
    \end{block}
    \vspace{0.3em}
    \begin{center}
        \begin{tabular}{|c|c|c|}
            \hline
            \textbf{层次} & \textbf{典型速度} & \textbf{差} \\
            \hline
            CPU寄存器 & 0.3ns & 1倍 \\
            \hline
            Cache（SRAM） & 1ns & $\sim$3倍 \\
            \hline
            主存（DRAM） & 50ns & $\sim$150倍 \\
            \hline
            磁盘 & 10ms & $\sim$3千万倍 \\
            \hline
        \end{tabular}
    \end{center}
    \vspace{0.3em}
    \begin{center}
        \textcolor{blue}{解决方案：在CPU和主存之间加一层小容量高速SRAM——Cache}
    \end{center}
\end{frame}

\begin{frame}{Cache的工作原理}
    \begin{block}{基本思想}
        \begin{enumerate}
            \item 将主存中\textbf{正在使用}的程序段和数据复制到Cache
            \item CPU访问存储器时，先检查Cache
            \item 若Cache命中（Hit），直接从Cache读取——快速
            \item 若Cache未命中（Miss），从主存读取，并复制到Cache
        \end{enumerate}
    \end{block}
    \vspace{0.3em}
    \begin{block}{关键概念}
        \begin{itemize}
            \item \textbf{命中率（Hit Rate）}：CPU访问在Cache中命中的比例
            \item \textbf{缺失率（Miss Rate）}：$1 - \text{命中率}$
            \item \textbf{平均访问时间}：
            \[
                T_{\text{avg}} = H \times T_{\text{Cache}} + (1-H) \times (T_{\text{Cache}} + T_{\text{主存}})
            \]
            \item $H$：命中率，希望越接近1越好
        \end{itemize}
    \end{block}
\end{frame}

\begin{frame}{局部性原理——Cache的理论基础}
    \begin{enumerate}
        \item \textbf{时间局部性}
        \begin{itemize}
            \item 刚被访问的单元不久后还可能被访问
            \item 例：循环变量、循环体指令反复执行
            \item $\Rightarrow$ 将数据保留在Cache中，下一次直接命中
        \end{itemize}
        \vspace{0.3em}
        \item \textbf{空间局部性}
        \begin{itemize}
            \item 刚被访问的单元附近也可能被访问
            \item 例：数组顺序访问、指令顺序执行
            \item $\Rightarrow$ 将相邻数据块（Cache行/块）一起调入
        \end{itemize}
    \end{enumerate}
    \vspace{0.5em}
    \begin{center}
        \textcolor{blue}{空间局部性 $\to$ 按块传输（Cache行大小通常为$2^n$字节）}
    \end{center}
\end{frame}

\begin{frame}{Cache的基本结构}
    \begin{block}{Cache行（Cache Line / Block）}
        每个Cache行包含：
        \begin{itemize}
            \item \textbf{有效位（Valid Bit）}：该行是否有效
            \item \textbf{标记（Tag）}：存储主存地址的高位，用于判断命中
            \item \textbf{数据块（Data Block）}：从主存复制的数据
        \end{itemize}
    \end{block}
    \vspace{0.3em}
    \begin{center}
        \begin{tabular}{|c|c|c|}
            \hline
            有效位 $V$ & 标记 Tag & 数据块 Data \\
            \hline
            1位 & $t$位 & $2^b$字节 \\
            \hline
        \end{tabular}
    \end{center}
    \vspace{0.3em}
    \begin{exampleblock}{地址划分}
        \[
            \underbrace{\text{Tag}}_{\text{标记}} \quad
            \underbrace{\text{Index/Set}}_{\text{索引}} \quad
            \underbrace{\text{Offset}}_{\text{块内偏移}}
        \]
        \begin{itemize}
            \item 标记：用于和Cache行中的Tag比较
            \item 索引：确定访问哪个Cache行/组
            \item 偏移：确定块内的具体字节
        \end{itemize}
    \end{exampleblock}
\end{frame}

\begin{frame}{Cache读操作流程}
    \begin{center}
        \begin{tabular}{c}
            CPU给出主存地址 \\
            $\downarrow$ \\
            解析为 Tag + Index + Offset \\
            $\downarrow$ \\
            根据Index找到Cache行 \\
            $\downarrow$ \\
            \begin{tabular}{c|c}
                有效位=1 且 Tag匹配？ & \\
                \hline
                是（命中） & 否（缺失） \\
                $\downarrow$ & $\downarrow$ \\
                返回数据块[Offset] & 从主存读取数据 \\
                                  & 写入Cache行并返回 \\
            \end{tabular}
        \end{tabular}
    \end{center}
    \vspace{0.3em}
    \begin{itemize}
        \item 写操作策略更加复杂（见下一页）
    \end{itemize}
\end{frame}

\begin{frame}{Cache写策略——写直达法}
    \begin{block}{写直达（Write-Through）}
        \begin{itemize}
            \item 写Cache的同时也写入主存
            \item 优点：主存数据始终保持一致
            \item 缺点：写操作速度慢（需要等主存完成）
            \item 改进：写缓冲（Write Buffer）—异步写入主存
        \end{itemize}
    \end{block}
    \vspace{0.5em}
    \begin{center}
        \begin{tabular}{c}
            CPU写数据 $\to$ Cache命中 \\
            $\downarrow$ \\
            同时写入Cache和主存 \\
        \end{tabular}
    \end{center}
\end{frame}

\begin{frame}{Cache写策略——写回法}
    \begin{block}{写回（Write-Back）}
        \begin{itemize}
            \item 写入时只修改Cache，不写主存
            \item 当Cache行被替换出Cache时，才写回主存
            \item 需要\textbf{脏位（Dirty Bit）}记录该行是否被修改
        \end{itemize}
    \end{block}
    \vspace{0.3em}
    \begin{center}
        \begin{tabular}{c}
            CPU写数据 $\to$ Cache命中 \\
            $\downarrow$ \\
            只修改Cache，置脏位=1 \\
            $\downarrow$ \\
            当该行被替换时，若脏位=1则写回主存
        \end{tabular}
    \end{center}
    \vspace{0.3em}
    \begin{itemize}
        \item 优点：减少主存访问次数
        \item 缺点：一致性维护复杂（多核时需\textbf{缓存一致性协议}）
    \end{itemize}
\end{frame}

\begin{frame}{写策略对比}
    \begin{center}
        \begin{tabular}{|c|c|c|}
            \hline
            \textbf{特性} & \textbf{写直达} & \textbf{写回} \\
            \hline
            写Cache时写主存 & 是 & 否 \\
            \hline
            主存一致性 & 好 & 需额外机制 \\
            \hline
            写速度 & 慢 & 快 \\
            \hline
            带宽占用 & 高 & 低 \\
            \hline
            是否需要脏位 & 否 & 是 \\
            \hline
            常用场景 & 简单系统 & 高性能系统 \\
            \hline
        \end{tabular}
    \end{center}
\end{frame}

\begin{frame}{写缺失时的处理策略}
    \begin{block}{写分配（Write-Allocate）}
        \begin{itemize}
            \item 写缺失时，先将主存块调入Cache，再写入
            \item 通常与\textbf{写回法}搭配使用
        \end{itemize}
    \end{block}
    \begin{block}{写不分配（No-Write-Allocate）}
        \begin{itemize}
            \item 写缺失时，直接写入主存，不调入Cache
            \item 通常与\textbf{写直达法}搭配使用
        \end{itemize}
    \end{block}
    \vspace{0.5em}
    \begin{center}
        \begin{tabular}{c|c}
            \textcolor{blue}{写回 + 写分配} & \textcolor{blue}{写直达 + 写不分配} \\
            \textcolor{blue}{（Write-Back + WA）} & \textcolor{blue}{（Write-Through + NW A）}
        \end{tabular}
    \end{center}
\end{frame}

\begin{frame}{本节总结}
    \begin{enumerate}
        \item Cache解决CPU-主存速度不匹配，基于局部性原理
        \item 时间局部性 $\to$ 保留数据，空间局部性 $\to$ 按块调入
        \item Cache行结构：有效位 + Tag + 数据块
        \item 写策略：写直达（同时写主存）vs 写回（替换时写回）
    \end{enumerate}
\end{frame}

\end{document}
